% Preface: the README's introduction, About the Course, Prerequisites and Written Papers sections,
% the course plan and the tools and hardware from info/README.md, and the build flow from
% hw/README.md and fw/README.md, as the front matter of a book.

\chapter{Preface}
\markboth{Preface}{Preface}

This book is the entry course to the \emph{FPGA-meets-MCU} pattern: build a hardware peripheral in
VHDL, write its driver in modern C++, then integrate the two across a real chip boundary. UART is
the simplest serial peripheral there is, chosen deliberately so that you meet the \emph{pattern},
the whole loop of hardware, driver and integration, before meeting a hard protocol. Harder
protocols apply the same loop later, in far greater depth; this is where its shape is learned.

The full stack, top (the application) to bottom (the wire), spans two chips joined by two serial
links. You write every layer of it except the SPI transport, which is provided:

\begin{booktable}{@{}p{14.6em}p{8.4em}L@{}}
  \thead{Layer} & \thead{Runs on} & \thead{Written by} \\
  \midrule
  \code{app::EchoNode} & ATmega328P (C++) & you \\
  The driver interface & ATmega328P (C++) & you \\
  \code{driver::uart::Uart} & ATmega328P (C++) & you \\
  The AVR SPI master, a \code{driver::transport::Interface} & ATmega328P (C++) & you \\
  \multicolumn{3}{@{}l}{\itshape The SPI link: four wires and a 5\,V to 3.3\,V level shifter (the
    control plane)} \\
  \code{spi_slave} and \code{spi_reg_bridge} & DE0-CV (VHDL) & provided, a black box \\
  The register bank & DE0-CV (VHDL) & you \\
  \code{uart_top}: baud generator, TX, RX, FIFOs & DE0-CV (VHDL) & you \\
  \multicolumn{3}{@{}l}{\itshape The UART link: \code{tx} and \code{rx} at 3.3\,V (the data
    plane)} \\
  A PC terminal on a USB-serial adapter & external & --- \\
\end{booktable}

Note the two serial links, and that they are \emph{not} the same protocol. The \textbf{control
plane} is SPI: the ATmega328P is the CPU, and it configures and drives the FPGA's UART peripheral
by reading and writing the peripheral's registers over SPI. The \textbf{data plane} is the UART
itself: the peripheral's own \code{tx} and \code{rx} lines carry bytes to and from the outside
world. There is no soft CPU anywhere; the processor is the real ATmega328P on an Arduino Nano.

\section*{Why drive a UART over SPI when the ATmega already has one?}
Because the ATmega's own UART is not the subject; the \emph{pattern} is. Here the ATmega plays host
CPU to a memory-mapped peripheral it does not contain, reaching it across a chip boundary. That is
the real situation with FPGA and ASIC peripherals, and it is what makes the skill transfer to the
next protocol. Building a UART this way is not the efficient way to get a UART; it is the way to
learn the boundary, with a protocol simple enough that nothing about the boundary is hidden by the
protocol's own complexity.

\section*{What the book covers}
The book covers the hardware/software boundary of an embedded system through the simplest serial
peripheral there is:
\begin{itemize}
  \item UART framing from the wire up: the idle-high line, the start bit, the data bits least
    significant first, optional parity, the stop bit or bits, and how a baud rate maps to a bit
    period.
  \item A baud-rate generator, a transmitter and an oversampling receiver in VHDL, each verified by
    a testbench.
  \item Register semantics over port semantics: poll-able \code{STATUS} bits and sticky
    \code{ERROR_FLAGS} bits built from single-cycle pulses; a write-triggered TX push; a read-then-pop
    RX path, and why a read that pops a FIFO would need the same commit-on-completion, abort-safe
    discipline as a write.
  \item The shared UART register protocol both halves implement against, the book's single source
    of truth, printed in full as Appendix~\ref{proto:ch}.
  \item A \code{driver::transport::Interface} seam under the C++ driver: \code{driver::uart::Uart}
    implements \code{driver::uart::Interface} over it, host-tested against a scripted
    \code{driver::transport::Stub} before any hardware exists.
  \item Freestanding C++ on the ATmega328P, and bring-up on the bench: two serial links, a level
    shifter, a bring-up ladder, and an echo application driven end to end.
\end{itemize}

\section*{What you are assumed to know}
The book is self-contained apart from the three foundations below, which it does \textbf{not}
teach again:
\begin{itemize}
  \item \textbf{Digital design with VHDL}: \code{entity} and \code{architecture}, processes,
    synchronizers, state machines, and running a provided self-checking testbench.
  \item \textbf{Modern embedded C++}: classes, interfaces, and the layered-driver idea.
  \item \textbf{Embedded C}: the ATmega toolchain (avr-gcc and avrdude) and register-level
    programming. It is needed only from \chapref{9}, when the driver reaches the real chip.
\end{itemize}

\subsection*{The provided transport}
The link between the MCU and the FPGA, a byte-level \code{spi_slave} and a transaction-level
\code{spi_reg_bridge}, is handed out as testbench-verified VHDL and used as a black box: you write
the peripheral behind it (the register bank and \code{uart_top}) and the driver above it, not the
bridge itself. Part~3 of the protocol (\secref{proto:part3}) documents the transport fully, so
nothing about it is magic, only out of scope. The box stays closed for one more course, since the
I2C course reuses the same control plane, and is opened in the one after: building
\code{spi_slave} and \code{spi_reg_bridge} from the wire up, clock-domain crossing and all, is the
whole subject of \emph{SPI: The MCU-FPGA Transport, from the Wire Up}.

\section*{How the book is organised}
The book is typeset from a ten-lecture course, and each chapter is one lecture. Its sections are the
lecture's appendices, in the same order and with the same subsections, so a chapter here and a
lecture directory in the repository can be read side by side. The lectures come in two parts, and so
does the book:

\begin{booktable}{@{}p{4.6em}p{4.4em}p{3.2em}L@{}}
  \thead{Chapter} & \thead{Lecture} & \thead{Side} & \thead{Topic} \\
  \midrule
  \multicolumn{4}{@{}l}{\itshape Part~I: The FPGA Peripheral} \\
  1 & \file{L01} & VHDL & The register package and the peripheral top \\
  2 & \file{L02} & VHDL & UART framing and the transmitter \\
  3 & \file{L03} & VHDL & The synchronizer and the receiver \\
  4 & \file{L04} & VHDL & The FIFO and the receive path \\
  5 & \file{L05} & VHDL & The register bank \\
  \midrule
  \multicolumn{4}{@{}l}{\itshape Part~II: The Driver and Integration} \\
  6 & \file{L06} & C++ & The driver's contracts \\
  7 & \file{L07} & C++ & The driver, built \\
  8 & \file{L08} & C++ & The driver, tested \\
  9 & \file{L09} & C++ & The real transport, and both toolchains \\
  10 & \file{L10} & Both & Integration and bring-up \\
\end{booktable}

The rhythm is five chapters of VHDL, four of C++, and one with both halves on the bench at once.
Part~I ends with a peripheral that passes its own system testbench in simulation; Part~II crosses
the wire to the ATmega and ends with both chips talking on the bench.

Every chapter opens with what is in it, works through the material, closes with a review of what
you should now be able to explain, and ends with its exercises. The appendices hold what every
chapter leans on. Appendix~\ref{proto:ch} is the UART register protocol, the contract both halves
are written against: read its Parts~2 and~3 (\secref{proto:part2} and \secref{proto:part3}) before
\chapref{1}, and Part~1 (\secref{proto:part1}) before \chapref{2}. Appendix~\ref{sol:ch} answers the
exercises that are not code, and the appendices after it hold the two written papers described
below.

\section*{What the exercises look like}
Roughly half the code the course produces lives in its exercises, and a provided self-checking
testbench or host suite is what says whether it works. The exercises are numbered within each
chapter, and each is labelled with its kind:

\begin{booktable}{@{}p{7.4em}L@{}}
  \thead{Kind} & \thead{What it asks of you} \\
  \midrule
  \kindtag[fillaccent]{VHDL} & Write a module from its section's specification, and check it with
    the provided testbench in GHDL. \\
  \kindtag[fillaccent4]{C++} & Write a header or a class from its specification, and check it with
    the provided host suite, \code{make test}. \\
  \kindtag[fillaccent2]{Reasoning} & Answer in words or numbers: explain why, trace what happens, or
    say what a testbench reports when you break a working design on purpose. \\
  \kindtag[fillaccent3]{Bench} & Work with the hardware in front of you: the Quartus flow, the
    wiring, and the bring-up ladder. \\
\end{booktable}

Most exercises mix the two halves of that table. A VHDL exercise typically starts with the module
and goes on to parts that ask what it does, what a deliberate mistake in it would do, and why the
specification is the way it is.

\textbf{The code is not answered in this book.} Each exercise names the testbench or the host suite
that checks it, and that check is the authority on your module: it is the executable form of the
specification, and a listing printed here would only be one answer of many. \textbf{Every other
part is answered}, in Appendix~\ref{sol:ch}: the reasoning, the numbers, the traces, and what a
testbench reports when you break a working design on purpose, verified by doing exactly that. Try
each part before you read its answer.

\section*{Two written papers, and what they are not}
Nothing in the course this book comes from is marked. Assessment is the exercises after every
chapter, the provided testbench or host suite that grades each one, and the questions in every
chapter's review.

Appendices~\ref{exam:ch:about} to \ref{ansb:ch:answers} hold two four-hour papers with worked
solutions, and they check something else: \textbf{what you can reconstruct on paper, with nothing
in front of you.} Eight questions each, covering all ten chapters, four of VHDL, three of C++ and
one with both halves on the bench, mixing theory with code you either trace and repair or write
from scratch, and with no GHDL and no \code{make test} in the room to tell you which it is.

\textbf{They exist purely so that you can test your own skills and knowledge after working through
the book. They gate nothing, they are not a qualification, and no part of the book requires
them.}

\textbf{Take one once you have finished the book}, after \chapref{10} and the bring-up ladder. Both
papers draw on all ten chapters and on both sides of the wire, so sitting one partway through
examines material you have not read yet.

\section*{Setting up}
Everything short of the bench runs in a Linux terminal: on Linux itself, or on Windows through WSL.
The versions below are the floor the build requires; the course is developed against newer ones.
\begin{itemize}
  \item \textbf{GHDL} 3.0 or newer, to analyze and simulate the VHDL. Everything is analyzed with
    \code{--std=93}.
  \item \textbf{g++} 7 or newer, for the host build and its tests, which are C++17.
  \item \textbf{avr-gcc} 5 or newer, with \textbf{avr-libc} and \textbf{avrdude}, to build and flash
    the ATmega328P. The AVR build is deliberately C++14, so that it also builds on the older
    compiler Microchip Studio ships.
  \item \textbf{clang-format}, for the style check.
  \item \textbf{Quartus Prime Lite} 20.1 or newer, which is the first with the Cyclone~V support the
    DE0-CV needs, to synthesize the peripheral and program the board. It is not needed until
    \chapref{9}, and it is a multi-gigabyte download, so install it before you get there.
  \item \textbf{A serial terminal}: \code{picocom}, \code{minicom} or \code{screen} on Linux, PuTTY
    on Windows, at 115200~8N1.
  \item \textbf{An editor}; Visual Studio Code is the one the course uses.
\end{itemize}

\noindent Then clone the companion repository, together with the test framework it keeps as a
submodule; in a clone made without \code{--recursive}, run \code{git submodule update --init} once:

\begin{shell}
git clone --recursive https://github.com/qrtech-academy/uart.git
\end{shell}

\noindent A few commands at its root cover everything the book asks of it:

\begin{shell}
make build           # Everything: the host C++ and its tests, then the VHDL.
make build-vhdl      # Analyze every module and run every testbench it can.
make build-cpp       # Build fw/ and run its host suites.
make format-check    # Check the formatting of the C++ and the VHDL.
\end{shell}

\noindent Code arrives chapter by chapter, and both builds are written for that. A testbench whose
modules you have not written yet is reported as skipped rather than failed, and so is a host suite
whose headers do not exist yet, so the build stays green while the design is still half-built.

All the VHDL lives in one directory, \file{hw/}: the modules you write, beside the provided
testbenches that check them, the provided transport and the provided register-map package. Nothing
is copied between chapters. A testbench is run from there in three GHDL steps, dependencies first;
for the transmitter of \chapref{2}:

\begin{shell}
cd hw
ghdl -a --std=93 uart_tx.vhd uart_tx_tb.vhd
ghdl -e --std=93 uart_tx_tb
ghdl -r --std=93 uart_tx_tb --assert-level=error
\end{shell}

\noindent \code{--assert-level=error} makes a failed check stop the simulation with a non-zero exit
code, rather than printing and running on to a misleading pass. Every testbench binds the module it
drives \textbf{positionally}, and so does every instantiation inside \code{uart_top}: the contract
is the order and the types of the ports, exactly as each section's interface table lists them, and
the names are yours.

The C++ lives in \file{fw/}: headers in \file{include/}, sources in \file{source/}, the host suites
in \file{test/}, and the ATmega328P port in \file{avr/}, which is cross-compiled and flashed with
\code{make -C fw/avr flash} rather than built on the host.

For the bench in \chapref{9} and \chapref{10} you need, per bench, a \textbf{Terasic DE0-CV}; an
\textbf{Arduino Nano} (any 5\,V, 16\,MHz ATmega328P Nano or clone); a \textbf{four-channel
bidirectional level shifter}, such as a BSS138-based board, for the four SPI lines; a \textbf{3.3\,V
USB-serial adapter}, such as a CP2102 or an FT232 set to 3.3\,V logic, as the far end of the data
plane; a breadboard and jumper wires; and the cables and the 12\,V supply for the DE0-CV's on-board
USB-Blaster. A logic analyzer is recommended but not required: any eight-channel 24\,MHz unit is
plenty at 1\,MHz \code{SCK} and 115200~baud, one exercise in \chapref{9} uses one and says so, and
nothing else depends on it.
